\section{Implementation}
\label{sec:implementation}

Chiyoda is implemented as a Python research toolkit. Scenario YAML files define
layouts, cohorts, hazards, responders, information parameters, intervention
policies, seeds, and variants. The scenario manager builds a simulation by
loading the layout, constructing agents, hazards, exits, responders,
navigation, spatial indexing, behavior state machines, and optional
intervention policies.

The core implementation lives in \texttt{chiyoda/core/simulation.py}. The
simulation loop records step telemetry, agent telemetry, cell-level maps,
exit flow, bottleneck dynamics, hazard state, gossip events, and intervention
events. Study execution lives in \texttt{chiyoda/studies/runner.py}; it
materializes variants, runs each seed, aggregates tables, and exports study
bundles. Analysis and figure export live in \texttt{chiyoda/analysis/}.

The information intervention layer lives in
\texttt{chiyoda/information/interventions.py}. It is deliberately small:
policies produce target points and messages, while the shared executor applies
messages to recipients and computes before/after entropy and accuracy. This
keeps policy design separate from belief mutation and telemetry.

The paper pipeline mirrors the code pipeline. Running
\texttt{python3 -m chiyoda.cli sweep scenarios/study\_information\_control.yaml}
emits tables and figures; \texttt{paper/scripts/gen\_stats.py} converts the
exported bundle into \texttt{stats.tex}; and the LaTeX build injects those
values into the abstract and evaluation.

\todo{Add an architecture figure once the first full study bundle has been generated.}
